\section{Dataset}
Neural network တွေကို သင်ပေးတဲ့နေရာမှာ သုံးတဲ့ dataset တွေအကြောင်း အရင်စပြောရအောင်။
အောက်မှာ ဥပမာ dataset တစ်ခုပေးထားပါတယ်။ ဒီ dataset မှာ data ၆ကြောင်းပါတယ်။
data တစ်ကြောင်းချင်းစီကို example သို့ sample တစ်ခုစီလို့ ခေါ်ဝေါ်ပါတယ်။

model ထဲကို input အနေနဲ့ထည့်မဲ့ $x$ မှာ ၁။ အလေးချိန်၊ ၂။ နှုတ်သီးအရှည်၊ ၃။ ဟောင်တတ်လား ဆိုတဲ့
အချက် သုံးခုရှိပါတယ်။ အဲတာကို အခေါ်အဝေါ်အားဖြင့် features တွေလို့ ခေါ်ပါတယ်။ မြန်မာလိုတိုက်ရိုက်ပြန်ရင်တော့ အရေအသွေးတွေ၊
ဂုဏ်သတ္တိတွေ လို့ ပြန်ဆိုရမှာပါ။ ဒါပေမဲ့ ကျွန်တော်တို့ ဒီဟာကို အင်္ဂလိပ်လို့ပဲ features လို့ပဲခေါ်သင့်ပါတယ်။

ဒါဆို ဒီ dataset မှာ feature ၃ခုရှိတယ်ပေါ့။ တစ်ကောင်ချင်းဆီရဲ့ကိုယ်အလေးချိန်၊ နှုတ်သီးအရှည်ရယ်၊ ဟောင်တတ်လား ဆိုပြီးတော့။
ဒီမှာ ဟောင်တတ်လားဆိုတာ ဟုတ်တယ်၊ မဟုတ်ဘူးပေါ့်။ ဟုတ်တယ်ဆိုရင် 1 မဟုတ်ဘူးဆိုရင်တော့ 0

model ကနေထုတ်ပေးမဲ့ $y$ ကို label သို့မဟုတ် class လို့ခေါ်တယ်။ ဒီမှာက ခွေးဖြစ်ရင် 1 ကြောင်ဆိုရင် 0 ပေါ့။ ဒီမှာမေးစရာက ဘာလို့ 0, 1
သုံးထားလဲပေါ့။ ကျွန်တော်တို့က neural network အခြေခံ အရိုးအရှင်းဆုံး model တွေကိုသုံးပြမှာ ဖြစ်လို့ပါ။ အခြေခံ model တွေဟာ စာလုံးတွေထက်
ဂဏန်းတွေကိုပဲ အသုံးပြုနိုင်တဲ့အတွက် ကျွန်တော်တို့က data ကို 0, 1 ဖြစ်အောင် ပြောင်းပေးရခြင်း ဖြစ်ပါတယ်။
function ပုံစံနဲ့ ရေးကြည့်မယ်ဆိုရင် ဒီအတိုင်း ရေးလို့ရပါတယ်။

$$
f(\text{အလေးချိန်}, \text{နှုတ်သီးအရှည်}, \text{ဟောင်တတ်လား}) = 0
$$

$$
f(\text{အလေးချိန်}, \text{နှုတ်သီးအရှည်}, \text{ဟောင်တတ်လား}) = 1
$$

\begin{table}[h]
\centering
\begin{tabular}{cccc}
\hline
$x_1$ (အလေးချိန် kg) & $x_2$ (နှုတ်သီးအရှည် cm) & $x_3$ (ဟောင်တတ်လား) & Class $y$ \\
\hline
12.0 & 8.5  & 1 & 1 \\
25.0 & 11.0 & 1 & 1 \\
8.5  & 7.0  & 1 & 1 \\
3.5  & 2.5  & 0 & 0 \\
4.2  & 3.0  & 0 & 0 \\
2.8  & 2.0  & 0 & 0 \\
\hline
\end{tabular}
\caption{Training Dataset: Dog vs Cat Classifier}
\label{tab:dataset}
\end{table}

\noindent ဒီလောက်ဆိုရင် dataset တွေရဲ့ အခြေခံ အခေါ်အဝေါ်တွေကို ရသွားပါပြီ။